\documentclass[11pt,a4paper]{article}
\usepackage[utf8]{inputenc}
\usepackage[T1]{fontenc}
\usepackage[margin=2.5cm]{geometry}
\usepackage{graphicx}
\usepackage{booktabs}
\usepackage{longtable}
\usepackage{hyperref}
\usepackage{natbib}
\usepackage{amsmath}
\usepackage{amssymb}
\usepackage{subcaption}
\captionsetup[subfigure]{labelformat=parens,labelsep=space}
\renewcommand{\thesubfigure}{\Alph{subfigure}}
\graphicspath{{results/figures_rev3/}{results/figures/}{results/phase6_scrna/}}

\title{Supplementary Material\\[0.4em]
\large Limited plasma proteomic coverage constrains genetic drug-target prioritization in colorectal cancer}
\author{Jiajie Zhou}
\date{}

\begin{document}
\maketitle

%------------------------------------------------------------------
\setcounter{figure}{0}
\renewcommand{\thefigure}{S\arabic{figure}}
\renewcommand{\figurename}{Supplementary Figure}
\subsection*{Supplementary Methods}

\paragraph{S1. Replication power approximation.}
The FinnGen directional-consistency analysis was not powered to replicate the
discovery associations, and a nominal $p<0.05$ in FinnGen should therefore not be
read as evidence against the prioritised genes. To quantify what the observed
replication rate implies, per-locus replication power was approximated as the
probability that the FinnGen inverse-variance-weighted test statistic reaches the
two-sided significance threshold, given a true effect equal to the discovery SMR
estimate:
\[
P_{\mathrm{repl}} \;=\; \Phi\!\left(\frac{|b_{\mathrm{discovery}}|}{SE_{\mathrm{FinnGen}}}-1.96\right)
\;+\; \Phi\!\left(-\frac{|b_{\mathrm{discovery}}|}{SE_{\mathrm{FinnGen}}}-1.96\right),
\]
where $b_{\mathrm{discovery}}$ is the discovery SMR effect and $SE_{\mathrm{FinnGen}}$
is the standard error of the inverse-variance-weighted replication estimate obtained
from the same cis instruments. Applied to the eight non-MHC Tier 1 loci, this gives a
mean expected replication proportion of 0.39 (per-locus range 0.05--0.86): PNKD 0.86,
GPBAR1 0.73, TMBIM1 0.59, COX14 0.43, CASC8 0.32, LAMC1 0.07, LIMA1 0.05 and
UTP23 0.05. The observed 1/8 nominal significance is therefore compatible with the
expected replication rate under a true-positive model, and the low FinnGen replication
rate is better attributed to the smaller replication sample than to absence of effect.
The calculation is reproduced by \texttt{scripts/supp\_replication\_power.py} from
\texttt{results/phase4\_finngen\_replication.csv}.

\subsection*{Supplementary Figures}

\begin{figure}[htbp]
\centering
\includegraphics[width=\textwidth]{FigS1a_manhattan.pdf}

\vspace{0.5em}

\includegraphics[width=0.72\textwidth]{FigS1b_qq_plot.pdf}
\caption{\textbf{Genome-wide SMR results.}
(Top) Manhattan plot of 15,582 gene-level SMR tests. Red line: Bonferroni threshold
($p < 3.21 \times 10^{-6}$). (Bottom) QQ plot with genomic inflation factor $\lambda = 1.600$.}
\label{fig:S1}
\end{figure}

\begin{figure}[htbp]
\centering
\includegraphics[width=\textwidth]{FigS2_tier1_regional_plots.pdf}
\caption{\textbf{Regional association plots for the six non-MHC Tier 1 genes.}
CRC GWAS $-\log_{10}p$ (blue) with SuSiE credible-set SNPs (red) and gene position (green dashed)
for UTP11, BMP2, PNKD, LAMC1, TMBIM1 and GPBAR1. Each locus resolved to single-SNP credible sets
(PIP = 1.0); panel titles give coloc PPH4.}
\label{fig:S23}
\end{figure}

\begin{figure}[htbp]
\centering
\includegraphics[width=\textwidth]{FigS3_susie_sensitivity.pdf}
\caption{\textbf{SuSiE fine-mapping sensitivity analysis.}
Heatmap of credible set counts across 108 parameter combinations for the six non-MHC
Tier 1 genes (65/108 successful; 43 failures concentrated in chr2:219Mb at $\pm$1,000 kb
windows).
All successful runs produced single-SNP credible sets (PIP = 1.0) across all conditions,
consistent with, but not by itself establishing, a single causal variant at these loci.}
\label{fig:S3}
\end{figure}

\begin{figure}[htbp]
\centering
\includegraphics[width=0.8\textwidth]{FigS4_smr_or_forest.pdf}
\caption{\textbf{Per-SD odds ratios for six non-MHC Tier 1 genes.}
Forest plot of SMR-derived per-standard-deviation odds ratios (OR) and 95\% confidence intervals for CRC risk.
BMP2 (OR = 1.38, 95\% CI 1.27--1.50) and TMBIM1 (OR = 1.35, 95\% CI 1.25--1.46) showed the largest risk
effects; GPBAR1 (OR = 0.78, 95\% CI 0.73--0.83) was the only protective gene.}
\label{fig:S2}
\end{figure}

\begin{figure}[htbp]
\centering
\includegraphics[width=0.8\textwidth]{FigS5_mr_sensitivity.pdf}
\caption{\textbf{HEIDI heterogeneity versus SMR significance.}
Scatter of $-\log_{10}(p_{\mathrm{HEIDI}})$ against $-\log_{10}(p_{\mathrm{SMR}})$ for the 75
Bonferroni-significant genes, colored by sensitivity tier. Dashed vertical line:
$p_{\mathrm{SMR}} = 3.21\times10^{-6}$ (Bonferroni threshold); dashed horizontal line:
$p_{\mathrm{HEIDI}} = 0.01$ (heterogeneity threshold).}
\label{fig:S6}
\end{figure}

\begin{figure}[htbp]
\centering
\includegraphics[width=0.6\textwidth]{FigS6_ukbppp_pqtl_coloc.pdf}
\caption{\textbf{UKB-PPP pQTL colocalization.}
Bar chart of coloc.abf PPH4 values for $n = 5$ UKB-PPP pQTL--GWAS colocalization tests
(TNF 0.2751, LTA 0.2605, TET2 0.0059, NCF2 $<$ 0.001, CDKN1A $<$ 0.001). All genes failed
coloc (PPH4 $<$ 0.8, dashed line).}
\label{fig:S4}
\end{figure}

\begin{figure}[htbp]
\centering
\includegraphics[width=0.6\textwidth]{FigS7_decode_pqtl_coloc.pdf}
\caption{\textbf{deCODE pQTL colocalization.}
Bar chart of coloc.abf PPH4 values for $n = 5$ deCODE pQTL--GWAS colocalization tests
(LIMA1 $<$ 0.001, CCM2 0.9514, STAT6 $<$ 0.001, TNFRSF1A $<$ 0.001, TNFRSF1B 0.0084). Only CCM2 passed
coloc (PPH4 = 0.9514).}
\label{fig:S5}
\end{figure}

\begin{figure}[htbp]
\centering
\includegraphics[width=0.8\textwidth]{FigS8_pqtl_mr_decode_forest.pdf}
\caption{\textbf{deCODE pQTL Mendelian randomization forest.}
Wald ratio (single top cis-pQTL) and inverse-variance-weighted (IVW) estimates with 95\% CI for
$n = 5$ deCODE genes with cis-pQTL instruments (LIMA1, CCM2, STAT6, TNFRSF1A, TNFRSF1B).
LIMA1 (Wald $p = 4.29\times10^{-7}$), CCM2 ($p = 6.28\times10^{-6}$) and STAT6
($p = 3.86\times10^{-5}$) were significant; TNFRSF1A ($p = 0.349$) and TNFRSF1B ($p = 0.614$)
were not.}
\label{fig:S20}
\end{figure}

\begin{figure}[htbp]
\centering
\includegraphics[width=0.7\textwidth]{FigS9_pqtl_mr_ukbppp_forest.pdf}
\caption{\textbf{UKB-PPP pQTL Mendelian randomization forest.}
Wald ratio and IVW estimates for $n = 4$ UKB-PPP genes with Olink protein measurements and
cis-pQTL instruments (CDKN1A, LTA, NCF2, TNF). CDKN1A (Wald $p = 0.020$, $b = +0.130$,
risk-increasing) and TNF ($p = 5.0\times10^{-4}$, $b = -0.295$, protective) were significant.}
\label{fig:S21}
\end{figure}

\begin{figure}[htbp]
\centering
\includegraphics[width=\textwidth]{FigS10_pqtl_mr_scatter.pdf}
\caption{\textbf{pQTL MR instrument scatter.}
SNP effect on protein level (exposure, pQTL $\beta$) versus SNP effect on CRC risk (outcome,
GWAS $\beta$) for the multi-instrument pQTL MR genes (CCM2, LIMA1, STAT6, CDKN1A, TNF); dashed
line = IVW causal estimate per gene.}
\label{fig:S24}
\end{figure}

\begin{figure}[htbp]
\centering
\includegraphics[width=\textwidth]{FigS11_pqtl_mr_loo.pdf}
\caption{\textbf{pQTL MR leave-one-out sensitivity.}
IVW causal estimate after sequentially dropping each instrument (triangle = estimate using
all instruments); panels show the five pQTL MR genes. Corner annotations give the largest
single-instrument shift ($|\Delta b|$) and the responsible instrument; for genes with more than
60 instruments the individual confidence intervals are omitted because the rows are not
resolvable at print size. All estimates remain within the same direction and magnitude,
indicating that no single instrument drives the result.}
\label{fig:S25}
\end{figure}

\begin{figure}[htbp]
\centering
\includegraphics[width=\textwidth]{FigS12_pqtl_mr_methods_forest.pdf}
\caption{\textbf{pQTL MR multi-method comparison.}
IVW, weighted median, weighted mode and MR-Egger slope estimates with 95\% CI for each of the five
pQTL MR genes. MR-Egger slope accounts for directional pleiotropy; qualitatively consistent
directions across methods support a robust causal estimate.}
\label{fig:S26}
\end{figure}

\begin{figure}[htbp]
\centering
\includegraphics[width=\textwidth]{FigS13_pqtl_mr_funnel_pleiotropy.pdf}
\caption{\textbf{pQTL MR funnel plot and pleiotropy diagnostics.}
Per-SNP Wald-ratio causal estimate versus instrument precision (1/SE) for the five pQTL MR genes;
symmetry around the IVW line indicates no directional pleiotropy. Annotations report MR-Egger
intercept and Cochran's Q. Statistically significant MR-Egger intercepts for STAT6
($p = 3.0\times10^{-12}$), CDKN1A ($p = 1.8\times10^{-4}$) and TNF ($p < 10^{-6}$) indicate
directional pleiotropy at these loci; high Cochran's Q indicates residual heterogeneity.
MR-PRESSO flagged 5 (STAT6), 10 (LIMA1), 15 (CCM2) and 22 (CDKN1A) outlier instruments and left
the direction of every estimate unchanged (Table~3).}
\label{fig:S27}
\end{figure}

\begin{figure}[htbp]
\centering
\includegraphics[width=\textwidth]{FigS14_celltype_dotplot.pdf}
\caption{\textbf{Cell-type specificity of SMR genes in CRC single-cell RNA-seq.}
Dotplot of the 16 most cell-type-specific SMR genes across 10 cell types in CRC single-cell RNA-seq
($n = 50,214$ cells; dot size = fraction expressing, color = mean normalized expression). MMP11 was
the only fibroblast-specific SMR gene (specificity = 5.78).}
\label{fig:S7}
\end{figure}

\begin{figure}[htbp]
\centering
\includegraphics[width=\textwidth]{FigS15_spatial_tier1_genes.pdf}
\caption{\textbf{Spatial expression of the six core Tier 1 candidates in CRC tissue.}
Visium spatial expression of UTP11, BMP2, PNKD, LAMC1, TMBIM1 and GPBAR1 across
$n = 3{,}493$ spots from one CRC section (GSE285505); colour scale is SCT-normalized
expression, scaled separately for each gene, and the percentage of spots with non-zero
expression is given under each gene name. Five of the six candidates were detected in
$11$--$57$\% of spots (UTP11 11.2\%, BMP2 12.5\%, PNKD 22.5\%, LAMC1 40.3\%,
TMBIM1 57.4\%), whereas GPBAR1 was detected in a single spot (0.03\%) and therefore
appears essentially blank. All six genes are shown so that tissue-level coverage can be
compared directly with the protein layer, in which only BMP2 had a matching plasma
measurement; TMBIM1 was additionally the only SMR gene with fibroblast-enriched
expression in the single-cell atlas (Supplementary Fig.~\ref{fig:S7}).}
\label{fig:S8}
\end{figure}

\begin{figure}[htbp]
\centering
\begin{subfigure}[b]{0.52\textwidth}
\centering
\includegraphics[width=\linewidth]{FigS16a_microenvironment_elbow_plot.pdf}
\caption{}
\end{subfigure}
\hfill
\begin{subfigure}[b]{0.46\textwidth}
\centering
\includegraphics[width=\linewidth]{FigS16b_me_zone_barplot.pdf}
\caption{}
\end{subfigure}
\\[0.6em]
\begin{subfigure}[b]{0.88\textwidth}
\centering
\includegraphics[width=\linewidth]{FigS16c_me_tier1_heatmap.pdf}
\caption{}
\end{subfigure}
\caption{\textbf{Spatial microenvironment zoning of CRC Visium data.}
(A) K-means within-cluster sum of squares for $k = 2$ to $k = 8$; the elbow at $k = 6$ defined
the final partition. (B) Composition of the six K-means-defined microenvironment (ME) zones
across 3{,}493 Visium spots; ME6 Stromal matrix was the largest zone (927 spots, 26.5\%),
followed by ME3 Mast cell zone (746, 21.4\%), ME5 T cell territory (613, 17.5\%), ME1
Myeloid-enriched (580, 16.6\%), ME2 Epithelial-dominant (518, 14.8\%) and ME4 B cell niches
(109, 3.1\%). (C) Expression of the 15 Tier 1 genes detected in Visium data across the six ME
zones (colour: mean SCT-normalised expression per spot; this panel reports
expression levels, not an enrichment score). Expression was not restricted to a
single zone: across the six core candidates the between-zone fold range in mean
expression was only 1.25--1.59, and UTP11, BMP2, LAMC1 and TMBIM1 each reached
their lowest mean expression in the stromal ME6 zone (PNKD: ME6 second lowest).
TMBIM1, which showed the strongest stromal association in the single-cell data,
was highest in the epithelial-dominant ME2 zone. GPBAR1 was detected in a single
epithelial ME2 spot (0.03\% of 3{,}493 spots).}
\label{fig:S14}
\end{figure}

\begin{figure}[htbp]
\centering
\includegraphics[width=0.9\textwidth]{FigS17_finngen_replication_forest.pdf}
\caption{\textbf{FinnGen replication Mendelian randomization forest.}
eQTL-IVW estimates of CRC risk in FinnGen R13 for $n = 8$ genes, coloured by direction
concordance with discovery SMR. Only LIMA1 achieved nominal significance (IVW $p = 0.0045$);
5 of 8 genes (LAMC1, TMBIM1, PNKD, GPBAR1, LIMA1) were direction-concordant.}
\label{fig:S22}
\end{figure}

\begin{figure}[htbp]
\centering
\includegraphics[width=0.8\textwidth]{FigS18_idg_druggability.pdf}
\caption{\textbf{IDG druggability grading of Tier 1 genes.}
Target Development Level (TDL) and tractability annotation (Pharos/DrugCentral). LAMC1 was the only
Tclin gene (approved drug ocriplasmin); GPBAR1 and BMP2 were Tchem; UTP11 was Tchem with antibody
tractability 0/3 (nuclear localization); TMBIM1 and PNKD were Tbio (no chemical probes or ligands).}
\label{fig:S18}
\end{figure}

\begin{figure}[htbp]
\centering
\includegraphics[width=0.7\textwidth]{FigS19_phewas_safety_scores.pdf}
\caption{\textbf{Composite PheWAS safety scores for Tier 1 genes.}
Composite scores for the six non-MHC Tier 1 genes from the Open Targets Platform
(2.0 $\times$ risk-category diversity + 1.5 $\times$ genetic evidence + 0.5 $\times$
high-risk count); UTP11, PNKD and LAMC1 scored 10.0. The score is dominated by the number of
genetic associations, so GPBAR1 (6.2) ranks below BMP2 (8.8) even though BMP2 has the narrower
organ-system risk spectrum (2 categories and 2 high-risk phenotypes, versus 2 categories and 8
high-risk phenotypes for GPBAR1). Table~1 reports the components separately and the
main text uses those.}
\label{fig:S19}
\end{figure}


%------------------------------------------------------------------
\clearpage
\setcounter{table}{0}
\renewcommand{\thetable}{S\arabic{table}}
\renewcommand{\tablename}{Supplementary Table}
\subsection*{Supplementary Tables}

\begin{table}[htbp]
\centering
\caption{Top 10 SMR associations for CRC (GCST90255675).}\label{tab:smr_top10}
\small
\begin{tabular}{rlllllll}
\toprule
  Rank & Gene & Chr & Top SNP & SMR $b$ & SMR $P$ & HEIDI $P$ & PPH4 \\
\midrule
  1 & CASC8 & 8 & rs28661658 & 0.236 & $3.36e-28$ & 0.0000 & --- \\
  2 & LAMC1 & 1 & rs11588675 & 0.134 & $7.14e-24$ & 0.6767 & 0.9294 \\
  3 & UTP23 & 8 & rs7829098 & -0.206 & $1.47e-23$ & 0.0000 & --- \\
  4 & LIMA1 & 12 & rs10876014 & 0.204 & $1.44e-19$ & 0.0000 & --- \\
  5 & COX14 & 12 & rs17124432 & 0.215 & $1.33e-15$ & 0.0000 & --- \\
  6 & PNKD & 2 & rs1473901 & 0.113 & $2.38e-15$ & 0.0007 & 0.9337 \\
  7 & ENSG:ENSG00000261338 & 2 & rs736730 & 0.060 & $3.43e-15$ & 0.0575 & --- \\
  8 & BMP2 & 20 & rs6140415 & 0.324 & $5.77e-15$ & --- & 0.9971 \\
  9 & GPBAR1 & 2 & rs11677953 & -0.251 & $1.75e-14$ & 0.1116 & 0.8132 \\
  10 & TMBIM1 & 2 & rs7559416 & 0.298 & $5.03e-14$ & 0.0016 & 0.8262 \\
\bottomrule
\end{tabular}

\tblnote{SMR tested 15582 genes; Bonferroni threshold $p < 3.21\times10^{-6}$.}
\end{table}

%% Table S2 is a full gene list (62 entries) typeset as a longtable that flows
%% across pages at 10 pt; give it a clean page start and keep the preceding float clear.
\clearpage
\small
\begin{longtable}{llllll}
\caption{Protein-coding genes (GENCODE v47 annotation) passing SMR Bonferroni significance ($n=62$).}\label{tab:protein_coding}\\
\toprule
  Gene & Chr & SMR $b$ & SMR $P$ & HEIDI $P$ & PPH4 \\
\midrule
\endfirsthead
\multicolumn{6}{l}{\textit{Table~\ref{tab:protein_coding} (continued)}} \\
\toprule
  Gene & Chr & SMR $b$ & SMR $P$ & HEIDI $P$ & PPH4 \\
\midrule
\endhead
\bottomrule
\endfoot
  LAMC1 & 1 & 0.134 & $7.14e-24$ & 0.6767 & 0.9294 \\
  UTP23 & 8 & -0.206 & $1.47e-23$ & 0.0000 & --- \\
  LIMA1 & 12 & 0.204 & $1.44e-19$ & 0.0000 & --- \\
  COX14 & 12 & 0.215 & $1.33e-15$ & 0.0000 & --- \\
  PNKD & 2 & 0.113 & $2.38e-15$ & 0.0007 & 0.9337 \\
  BMP2 & 20 & 0.324 & $5.77e-15$ & --- & 0.9971 \\
  GPBAR1 & 2 & -0.251 & $1.75e-14$ & 0.1116 & 0.8132 \\
  TMBIM1 & 2 & 0.298 & $5.03e-14$ & 0.0016 & 0.8262 \\
  UTP11 & 1 & 0.243 & $1.56e-13$ & 0.0467 & 0.9996 \\
  SMAD7 & 18 & 1.106 & $3.24e-13$ & 0.0000 & --- \\
  TMT1A & 12 & -0.206 & $3.36e-13$ & 0.0000 & --- \\
  SF3A3 & 1 & 0.148 & $1.11e-12$ & 0.0028 & --- \\
  CDKN1A & 6 & -0.288 & $1.51e-11$ & 0.1012 & --- \\
  ACTR1B & 2 & -0.199 & $1.77e-11$ & 0.3107 & --- \\
  POLD3 & 11 & -0.581 & $1.124e-10$ & 0.0000 & --- \\
  AAMP & 2 & 0.626 & $1.168e-10$ & 0.0155 & --- \\
  TNF & 6 & -0.169 & $1.310e-10$ & 0.1613 & --- \\
  B9D2 & 19 & -0.376 & $2.578e-10$ & 0.3501 & --- \\
  HLA-DRB5 & 6 & -0.054 & $6.720e-10$ & 0.0002 & --- \\
  STAB1 & 3 & 0.140 & $7.780e-10$ & 0.0004 & --- \\
  CABLES2 & 20 & -0.168 & $8.713e-10$ & 0.0000 & --- \\
  TMEM258 & 11 & -0.152 & $1.096e-09$ & 0.0001 & --- \\
  ATF1 & 12 & -0.146 & $1.968e-09$ & 0.0000 & --- \\
  STIMATE & 3 & -0.184 & $2.040e-09$ & 0.0023 & --- \\
  CERS5 & 12 & -0.070 & $2.116e-09$ & 0.0000 & --- \\
  CXCR1 & 2 & 0.137 & $4.493e-09$ & 0.0000 & --- \\
  CXCR2 & 2 & 0.165 & $6.285e-09$ & 0.0000 & --- \\
  LTA & 6 & -0.268 & $9.118e-09$ & 0.0386 & --- \\
  SMAD6 & 15 & 0.448 & $1.223e-08$ & 0.3097 & --- \\
  ASIC1 & 12 & -0.357 & $1.463e-08$ & 0.0000 & --- \\
  ARPC2 & 2 & 0.456 & $2.051e-08$ & 0.0025 & --- \\
  CCM2 & 7 & -0.085 & $2.322e-08$ & 0.1324 & --- \\
  ARPC5 & 1 & -0.076 & $3.218e-08$ & 0.2064 & --- \\
  HLA-E & 6 & 0.311 & $4.330e-08$ & 0.0015 & --- \\
  B3GNTL1 & 17 & -0.101 & $4.995e-08$ & 0.4413 & --- \\
  MZF1 & 19 & 0.157 & $5.422e-08$ & 0.0094 & --- \\
  TNNC1 & 3 & -0.399 & $5.997e-08$ & 0.0012 & --- \\
  METTL27 & 7 & -0.048 & $6.215e-08$ & 0.1033 & --- \\
  UQCC5 & 3 & -0.268 & $6.500e-08$ & 0.0001 & --- \\
  SMARCD1 & 12 & 0.357 & $7.903e-08$ & 0.0000 & --- \\
  CCDC97 & 19 & -0.683 & $1.047e-07$ & 0.2643 & --- \\
  SLC25A28 & 10 & -0.317 & $1.199e-07$ & 0.0000 & --- \\
  GDPGP1 & 15 & 0.186 & $1.430e-07$ & 0.0071 & --- \\
  TRIM4 & 7 & -0.056 & $1.800e-07$ & 0.0411 & --- \\
  TRIM28 & 19 & 0.323 & $2.178e-07$ & 0.0059 & --- \\
  LEMD3 & 12 & -0.093 & $2.386e-07$ & 0.4154 & --- \\
  ZMIZ1 & 10 & 0.292 & $2.677e-07$ & 0.1280 & --- \\
  RGL1 & 1 & -0.151 & $2.885e-07$ & 0.0583 & --- \\
  COX15 & 10 & -0.272 & $3.012e-07$ & 0.0000 & --- \\
  NCF2 & 1 & -0.267 & $4.092e-07$ & 0.0628 & --- \\
  EIF3H & 8 & -0.488 & $4.325e-07$ & 0.0000 & --- \\
  MMP11 & 22 & -0.350 & $5.084e-07$ & 0.0645 & --- \\
  PGAP3 & 17 & 0.111 & $6.926e-07$ & 0.1872 & --- \\
  STAT6 & 12 & -0.075 & $8.510e-07$ & 0.0000 & --- \\
  UBE2M & 19 & 0.305 & $1.271e-06$ & 0.0137 & --- \\
  CCHCR1 & 6 & 0.121 & $1.322e-06$ & 0.0000 & --- \\
  TET2 & 4 & -0.804 & $1.458e-06$ & 0.0074 & --- \\
  ZFP90 & 16 & 0.050 & $1.464e-06$ & 0.0006 & --- \\
  MYO9A & 15 & -0.193 & $1.725e-06$ & 0.0099 & --- \\
  ZFP57 & 6 & 0.041 & $2.211e-06$ & 0.0822 & --- \\
  CLDN4 & 7 & -0.331 & $2.781e-06$ & 0.0624 & --- \\
  ABHD12B & 14 & -0.113 & $3.104e-06$ & 0.4542 & --- \\
\end{longtable}

\tblnote{Genes are listed in order of decreasing SMR significance; SMR $b$ is the SMR effect size and HEIDI $P$ tests for heterogeneity between the eQTL and GWAS signals. PPH4 is given for the probe--gene pairs that passed Bayesian colocalization.}

\begin{table}[htbp]
\centering
\caption{SuSiE sensitivity across 18 parameter sets (SNP cap 200--500, $L$ 5--20, cis window $\pm$500--1000 kb). All 65 converged runs produced exclusively single-SNP credible sets (PIP=1.0). Because the number of credible sets returned tracked the user-specified $L$, this singleton pattern is interpreted in the main text as a property of the method and of the 1000 Genomes Phase 3 European reference panel rather than as evidence establishing a single causal variant per locus.}\label{tab:susie_sensitivity}
\small
\begin{tabular}{llllll}
\toprule
  Gene & Param combos & $n_{\text{CS}}$ range & $n_{\text{CS}}$ median & All single-SNP & Stability \\
\midrule
  BMP2 & 10 & 5--10 & 10 & Yes & Moderate \\
  GPBAR1 & 9 & 3--3 & 3 & Yes & Perfect \\
  LAMC1 & 15 & 5--20 & 10 & Yes & Moderate \\
  PNKD & 9 & 2--2 & 2 & Yes & Perfect \\
  TMBIM1 & 9 & 3--3 & 3 & Yes & Perfect \\
  UTP11 & 13 & 5--15 & 10 & Yes & Moderate \\
\bottomrule
\end{tabular}
\end{table}

\begin{table}[htbp]
\centering
\caption{UKB-PPP pQTL colocalization results. All 5 genes failed coloc (PPH4 $>$ 0.8). TNF and CDKN1A were significant in pQTL MR (Wald $p < 0.05$).}\label{tab:ukbppp_pqtl_coloc}
\small
\begin{tabular}{llllllll}
\toprule
  Gene & $n_{\text{SNPs}}$ & PPH4 & PPH3 & Coloc & $n_{\text{IV}}$ & Wald $b$ & Wald $p$ \\
\midrule
  TNF & 3504 & 0.2751 & 0.725 & Fail & --- & -0.295 & 0.0005 \\
  LTA & 3499 & 0.2605 & 0.739 & Fail & --- & -0.014 & 0.5060 \\
  TET2 & 809 & 0.0059 & 0.013 & Fail & --- & --- & --- \\
  NCF2 & 1135 & 0.0000 & 1.000 & Fail & --- & -0.071 & 0.3070 \\
  CDKN1A & 1307 & 0.0000 & 1.000 & Fail & --- & 0.130 & 0.0200 \\
\bottomrule
\end{tabular}
\end{table}

\begin{table}[htbp]
\centering
\caption{deCODE pQTL colocalization results. Only CCM2 passed coloc (PPH4 $>$ 0.8). The remaining 4 genes showed PPH3 $\approx$ 1.0, suggesting independent causal variants for pQTL and GWAS signals.}\label{tab:decode_pqtl_coloc}
\small
\begin{tabular}{llllllll}
\toprule
  Gene & Source & $n_{\text{SNPs}}$ & PPH4 & PPH3 & Coloc & Wald $b$ & Wald $p$ \\
\midrule
  STAT6 & deCODE & 5180 & 0.0000 & 1.000 & Fail & -0.029 & 0.0000 \\
  LIMA1 & deCODE & 5014 & 0.0002 & 1.000 & Fail & 0.858 & 0.0000 \\
  CCM2 & deCODE & 6850 & 0.9514 & 0.049 & Pass & -0.766 & 0.0000 \\
  TNFRSF1A & deCODE & 6859 & 0.0000 & 1.000 & Fail & -0.070 & 0.3487 \\
  TNFRSF1B & deCODE & 6462 & 0.0084 & 0.287 & Fail & -0.035 & 0.6139 \\
\bottomrule
\end{tabular}
\end{table}

\begin{table}[htbp]
\centering
\caption{Microenvironment zone characterization by K-means clustering (k=6) of module scores on 3493 Visium spots (GSE285505).}\label{tab:me_zones}
\small
\begin{tabular}{llllll}
\toprule
  Zone & N spots & \% & Cell type 1 & Cell type 2 & Cell type 3 \\
\midrule
  ME1 & 580 & 16.6\% & Myeloids (+1.59) & Stromal cells (0.13) & B cells (-0.12) \\
  ME2 & 518 & 14.8\% & Epithelial cells (+1.36) & Mast cells (-0.20) & T cells (-0.22) \\
  ME3 & 746 & 21.4\% & Mast cells (+1.42) & Stromal cells (0.13) & B cells (0.00) \\
  ME4 & 109 & 3.1\% & B cells (+3.66) & Stromal cells (0.86) & T cells (0.24) \\
  ME5 & 613 & 17.5\% & T cells (+1.51) & Stromal cells (-0.02) & B cells (-0.02) \\
  ME6 & 927 & 26.5\% & Stromal cells (+0.17) & B cells (-0.02) & Epithelial cells (-0.32) \\
\bottomrule
\end{tabular}
\end{table}

\clearpage
\input{manuscript/tables/TableS7_spatial_crosscorrelation.tex}
\begin{table}[htbp]
\centering
\caption{PheWAS safety screening of the six core non-MHC Tier 1 genes (Open Targets Genetics).}\label{tab:phewas_safety}
\small
\begin{tabular}{llllll}
\toprule
  Gene & Diseases & Genetic & Risk categories & Safety score & Tier \\
\midrule
  UTP11 & 25 & 19 & 4 & 10.0 & High alert \\
  PNKD & 25 & 9 & 4 & 10.0 & High alert \\
  LAMC1 & 25 & 20 & 3 & 10.0 & High alert \\
  TMBIM1 & 25 & 8 & 3 & 9.4 & High alert \\
  BMP2 & 25 & 20 & 2 & 8.8 & High alert \\
  GPBAR1 & 25 & 5 & 2 & 6.2 & High alert \\
\bottomrule
\end{tabular}

\tblnote{Safety score integrates disease count, genetic evidence strength, and risk category diversity.}
\end{table}

\clearpage
\input{manuscript/tables/TableS9_chen2024_gene_list.tex}

\end{document}
